\documentclass[11pt,a4paper,twocolumn]{article}
\usepackage{geometry}
\geometry{left=0.75in, right=0.75in, top=1in, bottom=1in, columnsep=0.25in}
\usepackage{times} % Standard academic font
\usepackage{hyperref}
\usepackage{authblk}
\usepackage{abstract}
\usepackage{cite}
\usepackage{booktabs}
\usepackage{multirow}
\usepackage{amsmath}
\usepackage{amssymb}
\usepackage{todonotes}

\title{\Large\textbf{Cross-Lingual Token Arbitrage: Optimizing Code Agent Context Windows via Local LLM Preprocessing}}
\author{\normalsize Utku \\ \textit{Software and Artificial Intelligence Engineer} \\ \textit{Istanbul, Türkiye}}
\date{}

\begin{document}

\twocolumn[
  \begin{@twocolumnfalse}
    \maketitle
    \begin{abstract}
AI-assisted coding agents are now bottlenecked by a single resource: the input token. Two pathologies of raw human input drive most of the overhead---a \textit{tokenization overhead} that charges $2$--$3\times$ more for non-English text, and \textit{structural entropy} in conversational prompts that triggers greedy retrieval and verbose decoding. Existing remedies act reactively, compressing already-bloated contexts or cascading on failure, without preventing bloat at its source.

We introduce a \textbf{pre-flight, edge-side prompt-rewriting middleware} that intervenes between the developer and the cloud agent. A localised \texttt{Llama 3.2} (3B) via \texttt{Ollama} performs three operations in a sub-second pass: cross-lingual translation into the cheaper English token space, structural rewriting into a compact \texttt{[CONTEXT]}/\texttt{[TASK]} format, and a regex-validated rewrite-with-fallback guarded by a $5\%$ token-budget threshold so the optimised payload is never larger than the original.

We evaluate on \textit{OMH-Polyglot}, the polyglot-specification arm of our \textit{OMH} (\textbf{O}ck-\textbf{M}ultilingual-\textbf{H}ard) benchmark family---$200$ instances whose technical specification is translated into Turkish, Arabic, Chinese, or code-switched mixtures, raising the per-prompt tokenization-overhead ratio from $0.96\times$ on the paired wrapper-only control \textit{OMH-Wrapped} to $2.05\times$ on average (max $6.15\times$). Across three commercial backends with three runs per condition, the middleware delivers prompt-side reductions of \textbf{$34$--$47\%$} and total-token reductions of $8.3\%$, $8.3\%$, and $18.8\%$ on \texttt{gpt-3.5-turbo}, \texttt{gpt-4o}, and \texttt{gemini-2.5-flash-lite} respectively. Accuracy is $99.5\%$ on \texttt{gpt-3.5-turbo} (unchanged), $99.5\%$ on \texttt{gpt-4o} ($+1.17$\,pp, within run-level noise), and $98.0\%$ on \texttt{gemini-2.5-flash-lite} ($+3.00$\,pp).



A \textit{names\_only} ablation isolates the deterministic function-name oracle from the SLM rewrite: the oracle alone \emph{decreases} accuracy by $18$--$20$\,pp on the two weaker backends (with a $3.8\times$ answer-token explosion on the lite model) and is neutral on \texttt{gpt-4o}, attributing the Gemini gain to the rewriter rather than to identifier salience. Against LLMLingua-2 \cite{pan2024llmlingua2} at matched compression rate, our middleware strictly Pareto-dominates on OckScore across all three backends, with LLMLingua-2 regressing accuracy by $-2.33$, $-22.17$, and $-63.00$\,pp respectively. At listed cloud prices, dollar-cost outcomes are asymmetric: $-12.4\%$ on Gemini, $-0.4\%$ on \texttt{gpt-4o}, and $+15.1\%$ on \texttt{gpt-3.5-turbo}, where the $3\times$ output/input price ratio inverts the token-side win.

We report all numerical contrasts as descriptive observations: OMH-Polyglot's $200$ rows are clustered into $20$ algorithm groups of $10$ style transforms with shared reference solutions, so the effective independent-trial count for accuracy is closer to $20$ than to $200$; an algorithm-cluster bootstrap and Holm-corrected significance pass are flagged as required follow-ups.
\vspace{1em}

\noindent \textbf{Keywords:} Token Arbitrage, Prompt Optimization, Code LLMs, Cost Efficiency, Retrieval-Augmented Generation, Local Preprocessing
\vspace{2em}

    \end{abstract}
  \end{@twocolumnfalse}
]

\section{Introduction}
AI-assisted Integrated Development Environments (IDEs) such as Cursor, GitHub Copilot, and Claude Code have become a default surface for professional software work, but their economic viability is now bottlenecked by a single resource: model's context window. Modern coding agents typically operate over repositories that exceed the model’s context window. To address this, they retrieve potentially relevant code fragments, prior interaction history, and auxiliary instructions before each generation step. Every chat round inflates the prompt with retrieved repository chunks, prior model turns, and unstructured user instructions. At deployment scale, even modest per-query overhead compounds into substantial financial and environmental cost \cite{du2026ockbench, peng2025coffe}.

We argue that a large and structurally avoidable fraction of this cost is paid \textit{before} the cloud model ever runs, and that it originates from two compounding pathologies of the raw human input:

\begin{enumerate}
    \item \textbf{Tokenization Overhead.} The subword tokenizers used by frontier LLMs (e.g., \texttt{cl100k\_base}) are statistically optimized for English. Equivalent semantic content in Turkish, Chinese, or Arabic can require $2\times$ to $3\times$ more tokens \cite{petrov2023language, teklehaymanot2025tokenization}, silently inflating both billing and effective context window for non-English-speaking developers.
    \item \textbf{Structural Entropy.} Real developer prompts are conversational, vague, and laden with social fillers. This ambiguity drives Retrieval-Augmented Generation (RAG) pipelines into ``greedy retrieval'' \cite{jimenez2024swebench}, where the agent over-fetches workspace context to compensate for an underspecified task, and induces verbose, redundant reasoning at decode time \cite{ahuja2025efficientxlang}.
\end{enumerate}


Existing remedies treat these symptoms reactively rather than at the source. Post-hoc token compressors such as LLMLingua \cite{jiang-etal-2023-llmlingua} and LongCodeZip \cite{shi2025longcodezip} prune an already-bloated context, paying a perplexity-computation overhead inside the developer's interactive loop. Cost-aware prompt-search frameworks such as EPiC \cite{taherkhani2025epic} and CAPO \cite{zehle2025capo} are offline, requiring expensive cloud-side iterations on the very API they aim to cheapen. Local--cloud cascading systems \cite{lu2026mccom, chen2024modelcascading} fall back to the large model with the original, unoptimized prompt the moment the small model fails. Crucially, none of these approaches proactively neutralize this tokenization overhead, and none enforce structural discipline on the input \textit{before} retrieval is triggered.

We therefore propose a different architectural primitive: a \textbf{pre-flight, edge-side prompt rewriter} placed between the developer and the cloud agent. A localized, parameter-efficient Small Language Model (\texttt{Llama 3.2} 3B \cite{meta2024llama32}) running through \texttt{Ollama} intercepts every raw query and, in a single sub-second pass, performs three operations: (i) \textit{Cross-Lingual Token Arbitrage}---translating non-English text into the cheaper English token space; (ii) \textit{Structural Rewriting} into a compact label-bracketed \texttt{[CONTEXT]} / \texttt{[TASK]} format (a \textit{Bi-Block} in the benchmark proxy, optionally extended to a Tri-Block with \texttt{[CONSTRAINTS]} in the IDE extension); and (iii) a regex-validated rewrite-with-fallback routine guarded by a $5\%$ token-budget threshold, so that the optimized payload is never larger than the original. Because the entire optimization runs on local hardware, the system incurs zero recursive cloud-API cost.

We evaluate this design on \textit{OMH-Polyglot} (Section~\ref{sec:omh-family}), the polyglot arm of our \textit{OMH} (\textbf{O}ck-\textbf{M}ultilingual-\textbf{H}ard) benchmark family, constructed specifically to stress both pathologies. Unlike \textit{OMH-Wrapped}---the wrapper-only member in which the algorithmic specification stays in English and only informal multilingual \emph{wrappers} (greetings, fillers, casual register) surround it---OMH-Polyglot translates the technical specification itself into Turkish, Arabic, Chinese, or genuine code-switched mixtures, raising the cl100k tokenization-overhead ratio from $0.96\times$ to $2.05\times$ on average. Across three commercial backends, evaluated with three runs per condition, the middleware reduces prompt tokens by a consistent \textbf{34--47\%} on \texttt{gpt-3.5-turbo}, \texttt{gpt-4o}, and \texttt{gemini-2.5-flash-lite}. Total-token savings are \textbf{$-8.3\%$, $-8.3\%$, and $-18.8\%$} respectively. Accuracy is unchanged on \texttt{gpt-3.5-turbo} ($99.5\%$), nominally $+1.17$\,pp on \texttt{gpt-4o} (within run-level noise), and $+3.00$\,pp on \texttt{gemini-2.5-flash-lite}. To attribute the Gemini gain we run a \textit{names\_only} ablation that strips out the SLM rewrite but keeps the deterministic function-name extraction; on the two weaker backends the oracle in isolation \emph{drops} accuracy by $18$--$20$\,pp (and on Gemini it triggers a $3.8\times$ answer-token explosion), confirming that the SLM rewrite is what holds and improves accuracy on multilingual inputs, not the identifier injection. We also report a finding that earlier pilot runs on \textit{OMH-Wrapped} (English spec with informal multilingual wrappers) had masked: answer-side compression is not universal---only \texttt{gemini-2.5-flash-lite} produces shorter completions under our rewrite; \texttt{gpt-4o} produces marginally longer ones ($+3.5\%$) and \texttt{gpt-3.5-turbo} produces substantially longer ones ($+42.3\%$), in both cases more than offset by the prompt-side savings. The previously reported ``Optimization Floor'' on \texttt{gpt-3.5-turbo} disappears on the cleaner benchmark, which we treat as a self-correction.

\paragraph{Contributions.} This paper makes four contributions:
\begin{itemize}
    \item We formalize \textit{Cross-Lingual Token Arbitrage} as an architectural mechanism, rather than the empirical observation it has been treated as in prior work \cite{ahuja2025efficientxlang}, and operationalize it as the first stage of a pre-flight middleware.
    \item We introduce a label-bracketed prompt schema (a \textit{Bi-Block} of \texttt{[CONTEXT]} and \texttt{[TASK]} for the benchmark setting; an optional Tri-Block extension with \texttt{[CONSTRAINTS]} for IDE deployments) and a regex-validated rewrite-with-fallback routine guarded by a $5\%$ token-budget threshold, so the optimizer can never inflate the cloud-billed payload.
    \item We release the \textit{OMH} (\textbf{O}ck-\textbf{M}ultilingual-\textbf{H}ard) benchmark family: \textit{OMH-Wrapped} (English spec with multilingual wrappers only) and \textit{OMH-Polyglot} (fully translated or code-switched specs), each $200$ instances over the same twenty hard cores and ten style indices. Polyglot raises the per-prompt tokenization-overhead ratio from $0.96\times$ to $2.05\times$ relative to Wrapped; we retain Wrapped as a tokenizer control (Section~\ref{sec:omh-family}) to show why wrapper-only prompts fail to surface the overhead we study.
    \item We provide an open implementation---an IDE extension and a headless evaluation proxy---and report cost--accuracy results on three commercial models, three runs each, showing consistent prompt-token reductions of $34$--$47\%$ and total-token reductions of $8.3$--$18.8\%$. We additionally observe accuracy that is unchanged or improved by up to $+3.00$\,pp depending on backend, and run a \textit{names\_only} ablation that isolates the deterministic function-name oracle from the SLM rewrite; the oracle in isolation \emph{decreases} accuracy by $18$--$20$\,pp on the two weaker backends and is neutral on \texttt{gpt-4o}, ruling it out as the explanation for the \textit{Ours} arm's accuracy results and attributing the Gemini gain to the rewriter. We further report a head-to-head comparison against the LLMLingua-2 \cite{pan2024llmlingua2} prompt compressor at matched compression rate (Section~\ref{sec:baselines}); our middleware strictly Pareto-dominates token-classification compression on OckScore across all three backends, and LLMLingua-2's failure pattern on the two weaker backends independently validates the structural-scaffold mechanism we identify in the \textit{names\_only} ablation. We disclose that (i) the $200$-row benchmark is clustered into $20$ algorithm groups, so accuracy contrasts still require an algorithm-cluster bootstrap and a multiple-comparisons-corrected significance pass before any claim about precise effect magnitude, and (ii) two cells of the original five-condition ablation matrix (\textit{regex\_strip}, \textit{llm\_no\_names}) remain pending (Section~\ref{sec:limitations}).
\end{itemize}

The remainder of the paper is organized as follows. Section~\ref{sec:related} situates our work against prompt compression, cost-aware prompt search, model cascading, and repository-level context management. Section~\ref{sec:methodology} details the middleware, including the Bi-Block / Tri-Block split between the benchmark proxy and the IDE extension and the regex-validated rewrite-with-fallback routine. Section~\ref{sec:setup} describes the OMH benchmarks and experimental protocol (Section~\ref{sec:omh-family} defines the family and naming), Section~\ref{sec:results} reports per-model results over three independent runs, Section~\ref{sec:conclusion} concludes, and Section~\ref{sec:limitations} enumerates the limitations of our approach.

\section{Related Work}
\label{sec:related}
To contextualize our contribution, we structure this review from the foundational use of LLMs in software engineering down to the specific works regarding token efficiency in agentic workflows.

\subsection{General Prompt Optimization and Cost Reduction}

The economic and computational overhead of processing extensive contexts in Large Language Models (LLMs) has catalyzed significant research into prompt optimization. Early automated prompt engineering frameworks, such as Automatic Prompt Engineer (APE) \cite{zhou2023large} and Automatic Prompt Optimization (APO) \cite{pryzant-etal-2023-automatic}, utilized LLMs to search for and refine instructions to maximize task accuracy. While effective for performance, these methods do not explicitly address the growing financial costs and latency associated with token bloat.

To directly mitigate inference costs, token-level prompt compression techniques have emerged. SelectiveContext \cite{li-etal-2023-compressing} utilizes a base language model to evaluate the self-information of lexical units, pruning redundant tokens based on their perplexity. Building on this, the LLMLingua framework \cite{jiang-etal-2023-llmlingua} employs a dynamic budget controller to selectively compress different prompt components, achieving high compression ratios for unstructured text. More recently, Evaluator Head-based Prompt Compression (EHPC) \cite{fei2025efficient} introduced a training-free approach that leverages specific attention heads in early transformer layers to skim and retain critical tokens during the pre-filling stage. Beyond static prompts, the AgentDiet framework \cite{xiao2026reducing} addresses agentic multi-turn bloat by employing a secondary, cost-efficient LLM to reactively summarize and prune expired trajectory history.

\subsubsection*{Limitations and the Need for Pre-Flight Middleware}

Despite these advancements, existing cost-optimization methods for code generation harbor significant operational gaps. First, optimization frameworks like EPiC still require API interactions with the primary LLM to evolve the prompts, while compression techniques like LongCodeZip impose high computational pre-processing overhead that disrupts seamless developer workflows \cite{shi2025longcodezip}. 

More critically, these approaches uniformly overlook the severe tokenization overhead imposed on non-English inputs by standard subword tokenizers. While recent empirical studies on multilingual models, such as EfficientXLang \cite{ahuja2025efficientxlang}, have demonstrated that cross-lingual reasoning (translating prompts to English before processing) can yield up to 40\% reductions in token usage without sacrificing accuracy, this phenomenon has primarily been treated as an analytical observation rather than an applied architectural mechanism. 

Furthermore, current code generation pipelines lack proactive \textit{structural enforcement}. They either rely on the human developer to consistently provide structurally disciplined prompts (a phenomenon highly susceptible to conversational noise and "user fault") or attempt to enforce constraints at decode-time, which still incurs the initial input token cost. 

Our work diverges fundamentally from these prior methodologies. Rather than relying on mathematically lossy token pruning or expending cloud API calls on iterative prompt evolution, we introduce a localized middleware architecture. By shifting the pre-processing burden to a parameter-efficient local LLM (e.g., Llama 3.2), our approach proactively executes \textbf{cross-lingual token arbitrage}—neutralizing tokenizer bias before dispatch—and strictly enforces a deterministic prompt structure. This ensures maximum semantic density and minimal token expenditure without incurring any recursive cloud API costs.

\subsection{Prompt Optimization and Token Economics in Code Generation}

The application of Large Language Models (LLMs) to software engineering has transitioned from foundational synthesis tools to highly autonomous, multi-turn agentic workflows. To maximize code correctness, state-of-the-art frameworks heavily rely on iterative execution feedback and self-reflection. However, these multi-turn interactions exponentially inflate context windows and cloud API costs, rendering them economically prohibitive for continuous IDE integration. This has driven the field to prioritize cost-efficiency specifically within code generation tasks through three primary avenues: Cost-Aware Prompt Search, Code-Specific Compression, and Model Cascading.

To minimize the financial burden of multi-turn interactions, researchers have introduced automated prompt search frameworks optimized for cost. Evolutionary Prompt Engineering for Code (EPiC) leverages a lightweight evolutionary algorithm to refine original prompts into high-quality versions with minimal LLM interactions \cite{taherkhani2025epic}. Similarly, Cost-Aware Prompt Optimization (CAPO) integrates budget-awareness directly into the prompt optimization loop using length penalties and statistical racing \cite{zehle2025capo}. While mathematically sound, these methods are fundamentally offline search algorithms that require extensive initial API interactions to evolve the prompt templates, making them unsuitable for intervening in real-time, unique developer queries.

To address the semantic fragmentation caused by general text compressors, domain-specific compression for code has emerged. LongCodeZip, for instance, employs coarse-grained compression that ranks and selects function-level chunks using conditional perplexity to preserve structural boundaries \cite{shi2025longcodezip}. However, it imposes high computational pre-processing overhead, relying on continuous perplexity calculations that disrupt the sub-second latency expected in interactive developer environments. 

The most recent paradigm in code generation economics is Model Cascading. Systems route simpler queries to smaller, cheaper language models and use self-generated tests as the cascade threshold to decide if the larger model is needed \cite{chen2024modelcascading}. Local-cloud frameworks, such as MCCom, explicitly cascade a local Small Language Model (SLM) with a cloud-based LLM, triggering the expensive cloud API only when the local SLM fails \cite{lu2026mccom}. Despite these advancements, Model Cascading remains reactive—if the smaller model fails, it passes the raw, unoptimized, and potentially bloated prompt directly to the expensive cloud model.

\subsubsection*{Structural Enforcement and Local Reasoning Middleware}

Current code optimization pipelines uniformly lack proactive \textit{structural enforcement} and overlook the severe tokenization overhead imposed on non-English inputs by standard subword tokenizers. Our work diverges fundamentally from reactive cascading and offline search by introducing a localized reasoning middleware. 

By shifting the pre-processing burden to a parameter-efficient local LLM (e.g., Llama 3.2), our architecture offloads the computational cost of prompt engineering entirely to the edge. Rather than attempting to solve the coding task itself, the local model actively reasons over the developer's intent, stripping away conversational fillers and executing \textbf{cross-lingual token arbitrage} to neutralize tokenizer bias before API dispatch. Most critically, it programmatically restructures the raw input into a deterministic format (e.g., \texttt{[CONTEXT]}, \texttt{[TASK]}, \texttt{[CONSTRAINTS]}). This pre-flight structural enforcement guarantees that the downstream cloud agent receives a hyper-dense, perfectly bounded specification, ensuring maximum semantic density and minimal token expenditure without incurring any recursive cloud API costs.

\subsection{Repository-Level Context Management for Code LLMs}

While function-level code generation has matured significantly, modern software engineering requires models to reason across entire codebases. Repository-level code generation necessitates understanding complex, long-range dependencies across disparate functions, classes, and modules \cite{liu2024graphcoder}. To achieve this, AI coding agents predominantly rely on Retrieval-Augmented Generation (RAG) to dynamically fetch workspace files into the prompt. Recent frameworks address this through various mechanisms: RepoCoder utilizes iterative retrieval and generation pipelines to repeatedly search the repository based on intermediate code drafts \cite{zhang-etal-2023-repocoder}, while systems like GraphCoder employ control-flow and data-dependence graphs to execute coarse-to-fine contextual retrieval \cite{liu2024graphcoder}.

\subsubsection*{The Criticality and Expense of Context Management}

Despite these architectural advancements, context management has emerged as the most critical and computationally expensive bottleneck in modern coding agents like Cursor and Claude Code. As agents accumulate multi-turn interaction histories and continuously ingest sprawling repository contexts, token usage exponentially inflates. In commercial deployments, every wasted token translates directly to increased financial cost and degraded inference latency. For instance, recent empirical observations highlight that efficient context engineering directly impacts an agent's economic viability; Claude Code, through improved context management, can utilize roughly 5.5x fewer tokens than competitors like Cursor for equivalent tasks, resulting in substantial API cost savings over time \cite{morphllm2026context}. Furthermore, unstructured and vague user prompts frequently trigger greedy retrieval mechanisms, causing RAG pipelines to over-fetch thousands of unnecessary tokens for minor codebase modifications.

\subsubsection*{Context Compression and Dependency Selection}

To mitigate the prohibitive costs of repository-level context bloat, recent literature has aggressively pivoted toward advanced context compression. Frameworks like SWEzze systematically distill and compress code contexts to their "minimal sufficient subsequences," removing distracting noise while preserving critical fix ingredients \cite{mechtaev2026compressing}. Evaluated on benchmarks like SWE-Bench, SWEzze demonstrated the ability to reduce total token budgets by over 50\% while simultaneously improving issue resolution success rates \cite{mechtaev2026compressing}.

A particularly vulnerable and critical aspect of repository-level generation is dependency selection. Providing an LLM with the exact, ground-truth dependency context is vital; missing these dependencies consistently leads to logic errors or redundant reimplementations \cite{lehai2026repository}. Conversely, when coding agents lack strict constraints or precise context, they frequently hallucinate or suggest vulnerable imports based on historical pre-training patterns rather than the actual repository architecture or local security requirements \cite{blackduck2026ossra}.

\subsubsection*{The Gap: Pre-Flight Context Discipline}

The existing paradigms for repository-level context management remain fundamentally reactive. Context compression models introduce their own inference-time latency overhead, and iterative retrieval pipelines do not prevent the initial RAG over-fetching triggered by chaotic human inputs. The literature distinctly lacks a mechanism that proactively structures the user's query \textit{before} the retrieval phase is executed. By utilizing a local reasoning middleware to explicitly define task constraints and enforce dependency boundaries prior to cloud dispatch, systems can dramatically narrow the scope of the downstream RAG search. This effectively neutralizes context bloat at its source, presenting a highly cost-efficient and structurally disciplined alternative to post-hoc compression.

\subsection{Benchmarking Code Generation: Quality vs. Token Economics}

Early evaluations of LLM-based code generation, such as HumanEval and MBPP, focused almost exclusively on functional correctness using metrics like Pass@k. As the field advanced to complex, repository-level reasoning, frameworks like SWE-bench continued this paradigm by evaluating models based primarily on their issue resolution rates \cite{jimenez2024swebench}. However, these traditional benchmarks largely ignore the computational and financial overhead required to achieve correctness, despite inference token consumption being a critical bottleneck for continuous IDE integration \cite{du2026ockbench}.

Recently, the literature has shifted toward evaluating "Sustainable Code Generation," recognizing that models achieving similar task accuracy can exhibit up to a 5x variance in token usage \cite{du2026ockbench, ali2026sustainable}. To address this, emerging benchmarks explicitly evaluate the trade-off between code quality and token efficiency. For instance, OckBench was introduced to jointly measure accuracy and token efficiency across reasoning and coding tasks, highlighting that token usage remains highly unoptimized in commercial LLMs \cite{du2026ockbench}. Similarly, the COFFE framework, introduced at FSE 2025, evaluates both the functional correctness and the execution efficiency of generated code using novel metrics such as \textit{Efficient@k} and stressful testing \cite{peng2025coffe}. In enterprise contexts, evaluations increasingly rely on cost-accuracy Pareto curves to assess the true operational viability of coding agents \cite{chen2024modelcascading}.

\subsubsection*{The Gap: From Post-Hoc Measurement to Proactive Optimization}

While these recent advancements provide the necessary metrics to quantify the token-cost-to-quality ratio, they serve primarily as post-hoc observational tools \cite{du2026ockbench, peng2025coffe}. They expose the underlying inefficiencies of current models but offer no integrated architectural mechanism to mitigate them prior to generation. Our research directly addresses this deficiency. Rather than merely benchmarking the inflated costs caused by raw, unstructured user prompts, our local middleware proactively optimizes the input architecture. By enforcing structural boundaries and executing cross-lingual token arbitrage, our approach artificially drives down the token consumption required to reach a correct solution, directly improving the system's standing on modern cost-efficiency benchmarks without sacrificing generation quality.

\section{Methodology}
\label{sec:methodology}
The proposed architecture functions as an autonomous, "pre-flight" middleware layer designed to intervene before the execution of the Retrieval-Augmented Generation (RAG) cycle in agentic IDEs. Unlike existing post-hoc compression techniques that prune tokens from an already bloated context \cite{jiang2023llmlingua}, our framework proactively optimizes the input architecture at the source. This is achieved by shifting the computational burden of prompt refinement to a localized, parameter-efficient Small Language Model (SLM), thereby neutralizing linguistic bias and structural entropy before they incur financial or computational costs in the cloud-billed environment.

\subsection{System Architecture and Middleware Orchestration}
The core of our approach is a localized preprocessing pipeline that manages the orchestration between the developer's raw input and the cloud-based frontier model (e.g., Claude 3.5 or GPT-4o). The system is implemented as a TypeScript-based gateway that intercepts queries within the IDE's internal request cycle \cite{peng2025coffe}. 

Our architecture utilizes a local \texttt{Ollama} instance running \texttt{Llama 3.2} (3B) \cite{meta2024llama32}. This specific model was selected due to its high reasoning density and optimized performance on edge hardware, enabling sub-second inference latency on standard development machines \cite{meta2024llama32}. The middleware enforces a strict separation of concerns through a local-cloud hybrid orchestration:
\begin{enumerate}
    \item \textbf{Local Reasoning Layer}: Responsible for cross-lingual translation, semantic stripping of "slop" (conversational fillers), and deterministic structural enforcement.
    \item \textbf{Cloud Execution Layer}: Reserved exclusively for high-complexity code synthesis based on the optimized, hyper-dense prompt received from the local layer.
\end{enumerate}

By intervening at the edge, the middleware prevents ``greedy retrieval'' patterns---where vague, unstructured prompts trigger the ingestion of thousands of irrelevant workspace tokens \cite{jimenez2024swebench}---and ensures that the downstream agent receives a deterministic specification. The empirical impact of this hybrid orchestration on three commercial backends is reported in Section~\ref{sec:results}; we anticipate the headline here only briefly: prompt tokens fall by $34$--$47\%$ uniformly across backends, and total tokens fall by $8.3$--$18.8\%$ without functional accuracy loss.

\subsection{Cross-Lingual Token Arbitrage}
A primary technical contribution of our framework is the execution of \textit{Cross-Lingual Token Arbitrage}. Modern subword tokenizers, such as \texttt{tiktoken} (CL100K\_base), utilize Byte-Pair Encoding (BPE) schemes that are statistically biased toward English linguistic patterns \cite{teklehaymanot2025tokenization}. Consequently, non-English text requires significantly more tokens to represent the same semantic information, imposing a measurable tokenization overhead on multilingual inputs \cite{petrov2023language}.

As demonstrated in Table \ref{tab:tokenization-overhead}, the tokenization penalty for non-English languages can exceed $3\times$ the cost of an equivalent English string. To neutralize this disparity, our middleware executes an algorithmic compression pass by programmatically translating non-English inputs into English at the local edge before cloud dispatch. 

\begin{table}[h]
\centering
\small
\begin{tabular}{lc}
\toprule
\textbf{Language} & \textbf{Relative Tokenization Cost} \\
\midrule
English & 1.00 \\
Spanish & 1.50 \\
Italian & 1.60 \\
Turkish & 2.16 \\
Chinese & 2.41 \\
Bulgarian & 2.60 \\
Arabic & 3.00 \\
\bottomrule
\end{tabular}
\caption{Relative tokenization cost compared to English for the \texttt{cl100k\_base} tokenizer.}
\label{tab:tokenization-overhead}
\end{table}

The values in Table \ref{tab:tokenization-overhead} highlight the economic and technical necessity of local translation. Foundational studies in tokenizer parity \cite{petrov2023language, teklehaymanot2025tokenization} confirm that these disparities lead to reduced effective context windows and inflated API latencies for non-English speakers. Our implementation utilizes the multilingual reasoning core of \texttt{Llama 3.2} \cite{meta2024llama32} to exploit the findings of \cite{ahuja2025efficientxlang}, which demonstrate that reasoning traces originally generated in English are more compact and semantically dense than their native counterparts in morphologically rich or low-resource languages. By using English as a token-efficient pivot language for semantic representation, the middleware achieves an immediate reduction in the input payload before structural optimization begins.

\subsection{Structural Rewriting: Bi-Block (Benchmark) and Tri-Block (IDE)}
\label{sec:structure}
A significant portion of token inefficiency in agentic IDEs originates from ``structural entropy'' in human-generated queries: conversational fillers, redundant social cues, and imprecise technical phrasing \cite{jiang2023llmlingua}. To mitigate this, the local SLM is asked to rewrite the raw user query into a fixed, label-bracketed format. We deploy two variants of this format that are deliberately distinct, since they target different operating contexts; we describe both honestly here so that the empirical claims in Section~\ref{sec:results}, which derive from the Bi-Block variant, are not confused with the larger Tri-Block design used in the IDE extension.

\paragraph{Bi-Block (used in all benchmark experiments).} The headless benchmark proxy (\texttt{benchmark-proxy/proxy.ts}) emits two labelled blocks per task:
\begin{itemize}
    \item \texttt{[CONTEXT]} contains a single fixed sentence: ``\textit{Python coding task; tests are authoritative.}'' This is a constant string in the proxy, \textbf{not} workspace-derived metadata; it functions as a generic preamble, not as a per-task grounding signal.
    \item \texttt{[TASK]} contains an SLM-produced one- or two-sentence English instruction stating what to implement, immediately followed by the original \texttt{assert} lines copied verbatim. The SLM is prompted to translate non-English text, strip conversational filler, and mention every required function name verbatim; it is explicitly forbidden from emitting code, \texttt{<solution>} envelopes, or section labels. Function names are additionally extracted from the \texttt{assert} lines by a deterministic regex and supplied to the SLM as a hint, so the cloud model receives correct identifiers even if the SLM's translation drifts.
\end{itemize}

This Bi-Block is the entire structural artifact produced for every result reported in Section~\ref{sec:results}. We make no claim that it isolates ``operational boundaries'' as a distinct block, because in OckBench-Coding the operational boundaries are encoded directly in the \texttt{assert} lines that we attach to \texttt{[TASK]}.

\paragraph{Tri-Block (IDE extension only, not measured here).} The interactive IDE extension (\texttt{src/promptOptimizer.ts}) targets a different setting in which workspace constraints are not given as executable tests, and therefore emits a third \texttt{[CONSTRAINTS]} block whose contents are whatever the SLM extracts from the raw query (style guides, naming conventions, performance targets, security notes that the user typed in natural language). The Tri-Block is therefore an architectural extension of the Bi-Block, not a separate verified-policy mechanism. We disclose this here because the empirical results in this paper bear only on the Bi-Block; any claim about the third block is a design proposal whose evaluation we defer.

\paragraph{Decoding.} In both variants the SLM is decoded greedily ($\texttt{temperature} = 0$, $\texttt{top\_p} = 0.9$, prediction budget $200$ tokens for the benchmark proxy and $400$ tokens for the IDE extension) with stop sequences that prevent the SLM from leaking code or section labels into the output.

By enforcing this structure prior to cloud dispatch, we align the prompt with the ``Sketch-of-Thought'' paradigm \cite{ma2025sketch}, which prioritises structural scaffolds over verbatim human text. We do not claim that the scaffold itself causes the answer-side savings reported in Section~\ref{sec:results}; that claim is examined empirically in Section~\ref{sec:answer-side}, where we show that answer-side behaviour is in fact backend-specific.

\subsection{Output Validation and Token-Budget Guard}
\label{sec:guard}
The risk of aggressive rewriting is that the SLM may emit malformed output (no \texttt{[CONTEXT]} / \texttt{[TASK]} markers, a Python code block instead of a specification, a filled \texttt{<solution>} envelope) or output that, after rewriting, is no smaller than the original prompt and therefore yields no benefit from being forwarded. We address both concerns with a deliberately small set of deterministic checks; we describe them in full here because the literal implementation is much simpler than ``closed-loop control'' would suggest.

\paragraph{Lightweight regex validation.} On every rewrite, the proxy applies a regex-based validator (\texttt{validateLight} in \texttt{benchmark-proxy/proxy.ts}; the IDE extension uses an analogous check) that triggers a regeneration if any of the following hold: (i) the output is empty or shorter than $10$ characters; (ii) the SLM accidentally emitted a Python code block matching the regex \verb|```\s*\w*[\s\S]*?\bdef\s+\w+\s*\(|; (iii) the SLM emitted a filled \verb|<solution>...</solution>| envelope; (iv) for the IDE Tri-Block variant, the required \texttt{[CONTEXT]} / \texttt{[TASK]} / \texttt{[CONSTRAINTS]} markers are not all present. If any condition fires, the proxy retries the rewrite with a repair prompt that names the violations; we allow up to two such repair attempts. We do not attempt to verify that the rewrite is semantically faithful to the original query: the rewrite is checked for syntactic legality only.

\paragraph{Token-budget guard.} After validation, the proxy applies a final guard that approximates token cost using a conservative character heuristic ($1$ token per non-ASCII character, $0.25$ tokens per ASCII character) and forwards the \emph{original} user prompt unchanged whenever the rewrite is not at least $5\%$ smaller than the raw input. This is a one-line safety net rather than a closed-loop control mechanism, but it is sufficient to ensure that the proxy can never inflate the cloud-billed payload relative to the raw baseline.

\paragraph{Failure mode.} If the rewrite repeatedly fails validation across all retries, the proxy logs the violations and forwards the raw user prompt to the cloud model. In the experiments of Section~\ref{sec:results} this fallback occurred on a small number of tasks per run; we report aggregate behaviour rather than separating fallback vs.\ successful-rewrite tasks, because the cloud-side cost of a fallback is already baked into the totals.

We acknowledge explicitly that this validation pipeline is much weaker than the multi-stage \emph{requirement extraction} / \emph{structural mapping} / \emph{constraint coverage} loop one might design for an industrial deployment. The implementation reported here is a pragmatic sanity check, and the gap between this and a full closed-loop verifier is a follow-up direction we list in Section~\ref{sec:limitations}.


\section{Experimental Setup}
\label{sec:setup}
To evaluate the efficacy of our cost-arbitrage middleware, we conduct a comparative analysis between raw cloud-based inference and our optimized local-cloud hybrid pipeline. While standard benchmarks often focus exclusively on accuracy, our evaluation prioritizes "Per-Token Intelligence"---the ability of a model to reach a correct solution with minimal token expenditure \cite{du2026ockbench}.

\subsection{The OMH benchmark family (Ock-Multilingual-Hard)}
\label{sec:omh-family}
We retain the OckBench-Coding \cite{du2026ockbench} \emph{harness}: each benchmark row is still a natural-language \texttt{problem} string followed by authoritative Python \texttt{assert}s, so functional grading and OckScore remain defined exactly as in OckBench. What we change is the \emph{content} of those rows. On the default OckBench-Coding split we used initially, the bundled coding tasks are comparatively easy for current commercial models: accuracy sits near ceiling and per-task token variance is modest, so the suite offers little headroom for stressing token-efficiency mechanisms or harder algorithmic reasoning.

We therefore introduce \textbf{OMH}, \textbf{O}ck-\textbf{M}ultilingual-\textbf{H}ard---a small benchmark \emph{family} of two $200$-row surfaces that plug into the same OckBench harness. Both surfaces share the \textbf{same twenty in-house algorithmic core problems} (reference implementations plus fixed \texttt{assert} suites) and the \textbf{same ten style indices} per core ($20\times 10=200$ rows each), generated by deterministic scripts checked into our artifact repository. The two members differ only in \emph{which parts of the prompt are allowed to be non-English}:
\begin{itemize}
    \item \textbf{\textit{OMH-Wrapped}} (\textit{Wrapper-only}): the technical specification for each core remains a short \textbf{English} sentence; each style index applies one of ten hand-written \emph{wrapper} templates (informal Turkish, Arabic, or Chinese openings; noisy or overly polite English; Turkish--English mix; etc.) around that fixed English spec. \textbf{Name:} \texttt{-Wrapped} signals that multilingual material wraps the spec without replacing it. This surface stresses conversational clutter and multilingual \emph{framing} while keeping an ``English floor'' on the actual task description; we use it mainly as a \textbf{tokenizer control} (Table~\ref{tab:tokenization-overhead-bench}).
    \item \textbf{\textit{OMH-Polyglot}} (\textit{Polyglot}): for each $(\text{core},\text{style})$ pair the \textbf{specification sentences themselves} are Turkish, Arabic, Simplified Chinese, or code-switched equivalents, sentence-aligned with the English authoring source; \texttt{assert} lines are never translated and stay ASCII so the grader is unchanged. \textbf{Name:} \texttt{-Polyglot} signals that the \emph{problem statement body} is genuinely multilingual. This is our \textbf{primary} surface for all accuracy, token, dollar, and baseline comparisons in Section~\ref{sec:results}.
\end{itemize}
Together, OMH-Wrapped and OMH-Polyglot isolate whether tokenization overhead and cross-lingual rewriting matter for \emph{wrapper noise} alone versus for \emph{fully translated specifications}---the tokenizer-ratio gap between them (Section~\ref{sec:benchmark-selection}) answers that question quantitatively.

\subsection{Ruling out alternative evaluation surfaces}
\label{sec:benchmark-selection}
Choosing the right evaluation surface is the single most consequential design decision in a token-economics paper, because almost every existing benchmark optimizes for the wrong axis. Besides the OMH family defined in Section~\ref{sec:omh-family}, we rule out three obvious off-the-shelf candidates explicitly:
\begin{itemize}
    \item \textbf{HumanEval, MBPP, SWE-Bench} \cite{jimenez2024swebench}: these score \texttt{pass@k} correctness without measuring the cost of correctness; tokens are not a first-class metric.
    \item \textbf{FLORES, XL-Sum and other translation benchmarks}: these measure translation fidelity rather than downstream code outcomes under multilingual input, so they cannot test whether neutralizing the tokenization overhead preserves task success.
    \item \textbf{Default OckBench-Coding} \cite{du2026ockbench}: while explicitly cost-aware, the default suite induces a ceiling effect for frontier models (near-perfect accuracy, low token variance) because the tasks are too easy to elicit divergent token usage.
\end{itemize}

\paragraph{Construction details for \textit{OMH-Wrapped}.}
The ten wrapper templates are materialized exactly as in our build script: Turkish casual; Arabic casual; Simplified Chinese casual; English with typos and chat register; overly polite English; Turkish formal-informal hybrid; Arabic formal; Chinese mixed-technical; stream-of-consciousness English; and Turkish--English code-mix. Every template interpolates the same per-core English specification string and leaves \texttt{assert}s untouched. A direct \texttt{cl100k\_base} measurement on the resulting JSONL gives a per-prompt tokenization-overhead ratio of only $0.96\times$ on the mean and $1.34\times$ at the $90$th percentile relative to an English-only floor, so OMH-Wrapped cannot stress-test compression of a fully non-English specification body.

\paragraph{Construction details for \textit{OMH-Polyglot}.}
The twenty cores are typical interview-style tasks (e.g., $0/1$ Knapsack, longest common subsequence, sliding-window maxima, topological sort, and Levenshtein distance). For each $(\text{core},\text{style})$ we \emph{replace} the English specification with sentence-aligned Turkish, Arabic, or Simplified Chinese authored alongside the English source. The ten style registers are: pure Turkish; pure Arabic (Modern Standard); pure Simplified Chinese; three pairwise sentence-alternating mixes (TR$\leftrightarrow$AR, TR$\leftrightarrow$ZH, AR$\leftrightarrow$ZH); a tri-language TR/AR/ZH rotation; and three native-narrative registers with embedded English CS jargon (one per base language). Only the \texttt{assert} lines remain ASCII so the OckBench grader is unchanged.

This construction raises the per-prompt tokenization-overhead ratio to $\mathbf{2.05\times}$ on the mean, $4.02\times$ at the $90$th percentile, and $6.15\times$ in the worst case (style \texttt{tr\_ar\_zh\_tri} on long specs). Table~\ref{tab:tokenization-overhead-bench} contrasts OMH-Wrapped with OMH-Polyglot and motivates evaluating the middleware on the latter.

\begin{table}[h]
\centering
\small
\begin{tabular}{lcc}
\toprule
\textbf{Benchmark} & \textbf{Mean overhead} & \textbf{p90 overhead} \\
\midrule
\textit{OMH-Wrapped} (wrapper-only control) & $0.96\times$ & $1.34\times$ \\
\textbf{\textit{OMH-Polyglot} (ours)}    & $\mathbf{2.05\times}$ & $\mathbf{4.02\times}$ \\
\bottomrule
\end{tabular}
\caption{Per-prompt tokenization-overhead ratio (\texttt{cl100k\_base} tokens of the multilingual prompt divided by an English-floor equivalent). \textit{OMH-Wrapped} adds informal multilingual wrappers around an English specification, so mean overhead stays near parity with the floor; \textit{OMH-Polyglot} translates the specification body and roughly doubles mean overhead (fourfold at the $90$th percentile).}
\label{tab:tokenization-overhead-bench}
\end{table}

We pair \textit{OMH-Wrapped} with \textit{OMH-Polyglot} when contrasting tokenizer behaviour (Table~\ref{tab:tokenization-overhead-bench}): the gap answers whether multilingual \emph{content} beyond pleasantries is required to surface tokenization overhead---by direct measurement, \textbf{yes}. All accuracy and cost results in Section~\ref{sec:results} are on \textit{OMH-Polyglot}; OMH-Wrapped appears only in this controlled tokenizer comparison.

\subsection{Inference Baselines and Implementation}
We evaluate our framework against three commercial cloud backends spanning the cost--capability frontier: OpenAI's \texttt{gpt-4o} and \texttt{gpt-3.5-turbo}, and Google's \texttt{gemini-2.5-flash-lite}. The local reasoning middleware uses a headless OpenAI-compatible proxy that runs \texttt{Llama 3.2} (3B) via an \texttt{Ollama} instance. The proxy is identical across backends; only the upstream cloud endpoint changes.

For the baseline (\textit{Raw}) runs, OckBench dispatches the multilingual prompt directly to the cloud API, with a fixed system prompt asking for a focused Python solution. For the optimized (\textit{Ours}) runs, the same OckBench harness routes its requests through our proxy on the same OpenAI-compatible endpoint; the proxy performs the local reasoning pass and forwards the rewritten payload, while leaving the system prompt unchanged. The system prompt is therefore identical across both arms, so any measured delta is attributable to the user-message rewrite alone. We run \textit{three} independent repeats per (model $\times$ pipeline) cell to characterize residual cross-region API non-determinism even at \texttt{temperature} $= 0$, and report mean and sample standard deviation throughout. All cloud calls use a $4096$-token output cap.

\subsection{Evaluation Metrics}
\label{sec:metrics}
For every (model $\times$ pipeline) cell we report the following:
\begin{itemize}
    \item \textbf{Accuracy}: percentage of tasks whose generated code passes every authoritative \texttt{assert}, computed as $\frac{1}{N}\sum_i \mathbb{1}[\text{task}_i \text{ passes}]$ with $N = 200$.
    \item \textbf{Token Count}: total, prompt-only, and answer-only token consumption summed across the $200$ tasks, together with the percent reduction relative to the same model's Raw mean.
    \item \textbf{OckScore}: a unified accuracy/cost metric introduced by \cite{du2026ockbench} that rewards accuracy and logarithmically penalises token waste. For direct comparability we adopt their definition,
    \begin{equation}
        \mathrm{OckScore} \;=\; \mathrm{Accuracy} \;-\; 10\cdot\log_{10}\!\Big(\tfrac{\bar{T}}{T_{\text{ref}}}\Big),
        \label{eq:ockscore}
    \end{equation}
    where $\mathrm{Accuracy}$ is in $[0, 100]$, $\bar{T}$ is the mean tokens-per-task on the run, and $T_{\text{ref}}$ is a fixed reference budget reported by the OckBench harness so that an efficient run is rewarded and a wasteful run is penalised. We reproduce the harness-reported OckScore values directly in Table~\ref{tab:main-results}; readers should interpret them comparatively (Ours vs.\ Raw within a model) rather than as an absolute quality score.
    \item \textbf{Wall-clock duration}: end-to-end seconds per $200$-task run, used to expose the latency overhead introduced by local rewriting (Section~\ref{sec:results-latency}).
\end{itemize}

\paragraph{Statistical scope.} We deliberately do not present confidence intervals or significance tests in this version of the paper, and we want to be explicit about why so reviewers can calibrate the strength of our claims:
\begin{itemize}
    \item \textbf{Cluster structure.} OMH-Polyglot is constructed as $20$ algorithmic core problems crossed with $10$ multilingual style transforms; the $10$ rows that share a core problem also share an identical reference solution and an identical \texttt{assert} suite. The effective number of independent algorithmic trials is therefore closer to $20$ than to $200$, and a defensible inferential analysis would use a cluster-level (per-algorithm) bootstrap rather than treating the $200$ rows as i.i.d.\ Bernoulli outcomes. We have not yet run that analysis.
    \item \textbf{Three-run determinism probe.} The three repeats per cell were performed at $\texttt{temperature}{=}0$ to characterise residual cross-region API non-determinism (Gemini was deterministic across all three runs; \texttt{gpt-3.5-turbo} varied by $<10^{-3}$; only one of the three \texttt{gpt-4o} Raw runs differed from the others, by a single task). Three deterministic-or-near-deterministic samples are a determinism probe, not a sample for population inference, and the run-level standard deviations we report should be interpreted as such rather than as input to a $t$-test.
    \item \textbf{Multiple comparisons.} If significance tests were added in a follow-up, the natural multiple-comparisons correction (e.g.\ Holm) would apply to $3$ backends $\times$ at least $2$ headline metrics; until that pass exists we present token-side numbers as descriptive contrasts and accuracy numbers as observations whose precise magnitude (though not whose direction or mechanism, see Section~\ref{sec:ablation}) is pending the cluster-level analysis.
\end{itemize}

We have run one of the originally-deferred single-arm ablations---\textit{names\_only}, which isolates the deterministic function-name oracle from the SLM rewrite---and report it in Section~\ref{sec:ablation}. The remaining two cells of the matrix (\textit{regex\_strip} and \textit{llm\_no\_names}) are still pending and are listed in Section~\ref{sec:limitations}.

\section{Results and Analysis}
\label{sec:results}
We evaluate the optimized pipeline (\textit{Ours}) against the raw cloud-call baseline (\textit{Raw}) on OMH-Polyglot, using three independent runs per cell. Table~\ref{tab:main-results} reports the run means; run-level sample standard deviations across the three repeats are uniformly small (accuracy $\le 1.15$\,pp; token columns $<0.5\%$ of the mean) and are listed in the supplementary material. As discussed in Section~\ref{sec:metrics}, those run-level numbers should be read as a determinism probe over a deterministic-or-near-deterministic backend, not as inferential standard errors over a sample of the algorithm population. Five descriptive findings emerge: (i) prompt-side cross-lingual arbitrage produces a large, within-backend token reduction across all three backends; (ii) answer-side compression is model-specific and partially contradicts the framing of prior work; (iii) the previously-reported ``Optimization Floor'' on \texttt{gpt-3.5-turbo} disappears on this cleaner benchmark; (iv) a \textit{names\_only} ablation (Section~\ref{sec:ablation}) shows that the deterministic function-name oracle in isolation \emph{harms} accuracy on the two weaker backends and is neutral on \texttt{gpt-4o}, attributing the Gemini accuracy gain to the SLM rewrite rather than to identifier injection; and (v) a head-to-head comparison against the LLMLingua-2 prompt compressor at matched compression rate (Section~\ref{sec:baselines}) shows our middleware strictly Pareto-dominates token-classification compression on OckScore across all three backends, with LLMLingua-2's failure pattern on the lite model independently validating the structural-scaffold mechanism.

The \textit{Ours} arm's structural-rewriting step is bundled with a deterministic oracle that extracts target function identifiers from the executable \texttt{assert} lines and places them at maximum salience inside the \texttt{[TASK]} block. We previously flagged this as a confound that left the accuracy deltas in Table~\ref{tab:main-results} potentially unattributable to the rewriter; the \textit{names\_only} ablation in Section~\ref{sec:ablation} now resolves that confound directly. The token-side claims in Table~\ref{tab:main-results} are unaffected by the oracle---prompt tokens are reduced by SLM rewriting regardless of identifier placement---and the accuracy claims survive the ablation: the oracle alone is insufficient (and on weaker models, harmful), so the rewriter is the active mechanism behind the accuracy column.

\begin{table*}[t]
\centering
\small
\setlength{\tabcolsep}{4pt}
\begin{tabular}{llccccc}
\toprule
\textbf{Cloud Agent} & \textbf{Pipeline} & \textbf{Acc.\ (\%)} & \textbf{Total Tokens} & \textbf{Prompt} & \textbf{Answer} & \textbf{OckScore $\uparrow$} \\
\midrule
\multirow{3}{*}{\texttt{gpt-3.5-turbo}}        & Raw         & $99.50$          & $94{,}338$                          & $53{,}713$          & $40{,}625$          & $99.04$ \\
                                               & LLMLingua-2 & $77.33$          & $85{,}026$                          & $42{,}314$          & $42{,}712$          & $76.91$ \\
                                               & Ours        & $\mathbf{99.50}$ & $\mathbf{86{,}474}$ ($-8.3\%$)      & $28{,}661$          & $57{,}812$          & $\mathbf{99.08}$ \\
\midrule
\multirow{3}{*}{\texttt{gpt-4o}}               & Raw         & $98.33$          & $139{,}085$                         & $43{,}565$          & $95{,}520$          & $97.66$ \\
                                               & LLMLingua-2 & $97.17$          & $124{,}790$                         & $37{,}281$          & $87{,}509$          & $96.56$ \\
                                               & Ours        & $\mathbf{99.50}$ & $\mathbf{127{,}594}$ ($-8.3\%$)     & $28{,}776$          & $98{,}819$          & $\mathbf{98.88}$ \\
\midrule
\multirow{3}{*}{\texttt{gemini-2.5-flash-lite}}& Raw         & $95.00$          & $116{,}653$                         & $44{,}918$          & $71{,}735$          & $94.43$ \\
                                               & LLMLingua-2 & $35.00$          & $243{,}193$                         & $37{,}932$          & $205{,}261$         & $33.85$ \\
                                               & Ours        & $\mathbf{98.00}$ & $\mathbf{94{,}725}$ ($-18.8\%$)     & $\mathbf{29{,}398}$ & $\mathbf{65{,}327}$ & $\mathbf{97.54}$ \\
\bottomrule
\end{tabular}
\caption{Raw cloud baseline, the LLMLingua-2 \cite{pan2024llmlingua2} prompt-compression baseline at matched compression rate ($r{=}0.6$, with the executable \texttt{assert} block detached and re-attached verbatim so the grader is never corrupted by lossy token pruning), and our optimized middleware on OMH-Polyglot ($200$ rows $=$ $20$ algorithm clusters $\times$ $10$ style transforms; three independent runs per cell). Reported values are run means; run-level sample standard deviations are small (accuracy $\le 1.15$\,pp; total tokens $<0.5\%$ of the mean for \textit{Ours} and the LLMLingua-2 baseline alike) and are omitted from the table for compactness---they are reported in the supplementary material and should be read as a determinism probe (Section~\ref{sec:metrics}) rather than as inferential standard errors. The percentage in parentheses is the relative change in total tokens versus the same model's Raw mean; per-column relative changes for Prompt and Answer are discussed in Sections~\ref{sec:results} and \ref{sec:baselines}. OckScore is computed via Equation~\ref{eq:ockscore}. Across all three backends \textit{Ours} strictly Pareto-dominates LLMLingua-2 on the OckScore axis ($99.08 > 76.91$ on \texttt{gpt-3.5-turbo}, $98.88 > 96.56$ on \texttt{gpt-4o}, $97.54 > 33.85$ on \texttt{gemini-2.5-flash-lite}).}
\label{tab:main-results}
\end{table*}

\subsection{Prompt-Side Cross-Lingual Arbitrage is the Robust Effect}
The prompt-token column in Table~\ref{tab:main-results} is the strongest and most internally consistent finding in our study: \textit{within each backend}, the optimized pipeline reduces that backend's own prompt-token count by between $34$\% and $47$\%, with run-to-run variation under $20$ tokens out of nearly $30{,}000$ on the optimized arm. Each backend uses a different tokenizer (\texttt{cl100k\_base} for \texttt{gpt-3.5-turbo}, \texttt{o200k\_base} for \texttt{gpt-4o}, and Gemini's SentencePiece-derived tokenizer for \texttt{gemini-2.5-flash-lite}), and the absolute token counts these tokenizers report are not commensurate---the same byte string maps to different counts under each scheme, so the post-rewrite numbers $28{,}661$ / $28{,}776$ / $29{,}398$ should not be read as a cross-tokenizer ``convergence'' to a single underlying length. The defensible reading is that \emph{each tokenizer}, applied to its own backend's prompts, separately exhibits a large reduction (Ours vs.\ Raw within a row of Table~\ref{tab:main-results}), and that reduction holds across three otherwise-disparate commercial APIs. We treat this within-backend reduction, not cross-tokenizer length equivalence, as the prompt-side claim of this paper.

\subsection{Answer-Side Behaviour is Model-Specific}
\label{sec:answer-side}
A more nuanced---and to us, more interesting---finding concerns the \textbf{answer} token column. Only \texttt{gemini-2.5-flash-lite} produces shorter completions under our rewrite ($-8.9\%$). Both OpenAI models produce \textit{longer} completions: $+3.5\%$ on \texttt{gpt-4o} and $+42.3\%$ on \texttt{gpt-3.5-turbo}. We deliberately do \textit{not} claim the rewrite ``cures the overthinking tax'' on these models. The most plausible explanation is that the structurally clean Bi-Block prompt, by removing conversational ambiguity, gives the cloud model less reason to hedge and more reason to fully elaborate the implementation. We also note that on \texttt{gpt-4o} the apparent $+1.17$\,pp accuracy delta corresponds to a single task swinging in one of the three Raw repeats; this is well within the run-level standard deviation ($1.15$\,pp) and the \textit{names\_only} ablation (Section~\ref{sec:ablation}) confirms that gpt-4o is essentially flat across all three arms ($98.33$ Raw, $98.50$ \textit{names\_only}, $99.50$ \textit{full}), so we do \textbf{not} claim the gpt-4o accuracy gain as a result. Section~\ref{sec:ablation} also surfaces a striking answer-side counterpart: under \textit{names\_only} on Gemini, mean answer length explodes to $1{,}353$ tokens---more than $4\times$ the \textit{full}-arm length of $327$---direct evidence that structural scaffolding is what \emph{prevents} the lite model from generating runaway output, not merely a token-saving by-product. Crucially, the prompt-side savings are large enough to absorb the answer-side increase in every case, so net total tokens still fall by $8.3$--$18.8\%$. The honest synthesis is that our middleware reliably compresses the \emph{input} surface; it shapes the \emph{output} surface only on certain backends, and on the lightest backend it specifically prevents destabilisation that the alternative arms make visible.

\subsection{The Optimization Floor Was a Benchmark Artifact}
\label{sec:floor}
Earlier preliminary experiments on \textit{OMH-Wrapped} (the wrapper-only hard benchmark) suggested an ``Optimization Floor'' for \texttt{gpt-3.5-turbo}, where compression was accompanied by a small accuracy regression. On the cleaner OMH-Polyglot benchmark, with three independent runs and an identical system-prompt policy across both arms, no such regression appears: \texttt{gpt-3.5-turbo} accuracy is exactly $99.50\%$ in every Raw and every Ours run. We interpret the previously reported floor as a benchmark and harness artifact rather than a model-architecture finding, and have updated the framing accordingly. The \textit{names\_only} ablation in Section~\ref{sec:ablation} sharpens this further: when we strip out the SLM rewrite but keep the multilingual body, \texttt{gpt-3.5-turbo} drops from $99.50\%$ to $79.00\%$, a $20.5$\,pp collapse. The ``floor'' was never a property of the model under compression; it is the symptom of a multilingual prompt being forwarded to a weaker backend without translation, and the rewriter is what removes it.

\subsection{Gemini's Token-Side Win and the Rewriter as the Active Mechanism}
\label{sec:gemini}
\texttt{gemini-2.5-flash-lite} is the only backend that simultaneously reduces prompt tokens ($-34.5\%$), reduces answer tokens ($-8.9\%$), and shows an accuracy gain ($+3.00$~pp, $95.00 \rightarrow 98.00$, six previously-failing rows now passing). Token-side this yields the largest OckScore improvement in Table~\ref{tab:main-results} ($94.43 \rightarrow 97.54$). The accuracy gain has a known potential confound---the \textit{Ours} arm bundles structural rewriting with a deterministic function-name oracle---which we resolve directly via the \textit{names\_only} ablation in Section~\ref{sec:ablation}: the oracle in isolation \emph{decreases} Gemini accuracy by $18.50$\,pp ($95.00 \rightarrow 76.50$) and induces a $3.8\times$ answer-token expansion (mean $1{,}353$ tokens vs.\ $359$ in Raw and $327$ in Ours). The salience-only hypothesis is therefore falsified by direct measurement; the $+3.00$\,pp gain in Table~\ref{tab:main-results} is attributable to the SLM rewrite (the only mechanism present in \textit{full} and absent in \textit{names\_only}), not to identifier injection. The cluster-level caveat from Section~\ref{sec:metrics} still applies to the precise magnitude of the gain (six row-flips concentrated in $\le 6$ algorithm clusters), but the attribution to the rewriting mechanism is now unambiguous. The token-side ($-18.8\%$ total, $-12.4\%$ in dollars) remains the strongest single-backend result in the paper.

\subsection{Ablation: Isolating the Function-Name Oracle}
\label{sec:ablation}
To attribute the accuracy column of Table~\ref{tab:main-results} to a specific mechanism, we run a third arm, \textit{names\_only}, in which the SLM is bypassed entirely and the regex-extracted function name is prepended to the raw multilingual prompt as a one-line English oracle. This isolates the effect of identifier salience in the absence of any English rewriting, structural scaffolding, or filler stripping. We run three independent repeats per (model $\times$ arm) cell at \texttt{temperature}\,$=0$, identical in all other respects to the runs reported in Table~\ref{tab:main-results}. Results are summarised in Table~\ref{tab:ablation}.

\begin{table*}[t]
\centering
\small
\setlength{\tabcolsep}{4pt}
\begin{tabular}{llcccc}
\toprule
\textbf{Cloud Agent} & \textbf{Arm} & \textbf{Acc.\ (\%)} & \textbf{Total Tokens} & \textbf{Prompt} & \textbf{Answer} \\
\midrule
\multirow{3}{*}{\texttt{gpt-3.5-turbo}}        & \textit{raw}         & $99.50$           & $94{,}338$           & $53{,}713$           & $40{,}625$ \\
                                               & \textit{names\_only} & $79.00 \pm 1.32$  & $100{,}960$          & $60{,}093$           & $40{,}867$ \\
                                               & \textit{full}        & $\mathbf{99.50}$  & $\mathbf{86{,}474}$  & $\mathbf{28{,}661}$  & $57{,}812$ \\
\midrule
\multirow{3}{*}{\texttt{gpt-4o}}               & \textit{raw}         & $98.33$           & $139{,}085$          & $43{,}565$           & $95{,}520$ \\
                                               & \textit{names\_only} & $98.50 \pm 1.32$  & $131{,}336$          & $49{,}212$           & $82{,}124$ \\
                                               & \textit{full}        & $\mathbf{99.50}$  & $\mathbf{127{,}594}$ & $\mathbf{28{,}776}$  & $98{,}819$ \\
\midrule
\multirow{3}{*}{\texttt{gemini-2.5-flash-lite}}& \textit{raw}         & $95.00$           & $116{,}653$          & $44{,}918$           & $71{,}735$ \\
                                               & \textit{names\_only} & $76.50 \pm 0.00$  & $321{,}654$          & $51{,}078$           & $\mathit{270{,}576}$ \\
                                               & \textit{full}        & $\mathbf{98.00}$  & $\mathbf{94{,}725}$  & $\mathbf{29{,}398}$  & $65{,}327$ \\
\bottomrule
\end{tabular}
\caption{Ablation isolating the function-name oracle. \textit{names\_only} prepends the regex-extracted function identifier to the raw multilingual prompt and skips the SLM rewrite; \textit{raw} and \textit{full} are reproduced from Table~\ref{tab:main-results} for direct comparison. Three runs per cell; $\pm$ values are the run-level standard deviation. Bold marks the best value within each model. Italic flags Gemini's $3.8\times$ answer-token expansion under \textit{names\_only} ($1{,}353$ tokens/answer vs.\ $359$ in raw and $327$ in full), evidence that the lite model destabilises into runaway generation when given a salient identifier without semantic context. Each \textit{names\_only} cell shows a $\sim$$6{,}000$-token jump in prompt tokens over its \textit{raw} row, exactly consistent with the $\sim$$30$-token oracle prefix on $200$ tasks, confirming the proxy is in the loop.}
\label{tab:ablation}
\end{table*}

The ablation cleanly separates the two hypotheses raised in Section~\ref{sec:gemini}.

\paragraph{The oracle alone is insufficient on weaker backends, and harmful.} On \texttt{gpt-3.5-turbo} the oracle drops accuracy by $20.50$\,pp ($99.50 \rightarrow 79.00$); on \texttt{gemini-2.5-flash-lite} by $18.50$\,pp ($95.00 \rightarrow 76.50$). The mechanism is consistent with under-attention to the multilingual body once a salient English instruction is prepended: the model treats the prefix as an authoritative spec and partially ignores the foreign-language text where edge cases and return semantics are encoded. On \texttt{gpt-4o} the oracle is essentially neutral ($+0.17$\,pp, well within the $\pm 1.32$\,pp run-level standard deviation), consistent with a frontier model strong enough to recover the spec from the multilingual body without help.

\paragraph{The rewriter, not the oracle, is the active accuracy mechanism.} If the function-name oracle were the source of the \textit{full} arm's accuracy delta, \textit{names\_only} should approximate \textit{full} accuracy. It does not: \textit{full} exceeds \textit{names\_only} by $20.50$\,pp on \texttt{gpt-3.5-turbo}, by $1.00$\,pp on \texttt{gpt-4o}, and by $21.50$\,pp on \texttt{gemini-2.5-flash-lite}. The salience-only hypothesis floated for Gemini in Section~\ref{sec:gemini} is therefore directly falsified, and the $+3.00$\,pp Gemini gain in Table~\ref{tab:main-results} is attributable to the SLM rewrite---the only mechanism present in \textit{full} and absent from \textit{names\_only}.

\paragraph{Gemini destabilises into runaway generation when given salience without scaffold.} Under \textit{names\_only}, Gemini's mean answer length explodes from $359$ tokens (\textit{raw}) and $327$ tokens (\textit{full}) to $1{,}353$ tokens---a $3.8\times$ expansion. With our $2048$-token output cap, a non-trivial fraction of these answers are likely cap-truncated, which contributes some share of the $-18.50$\,pp accuracy loss; the qualitative picture, even allowing for this, is that the lite model fails to terminate cleanly when given a salient identifier in the absence of English semantic scaffolding. This is direct empirical support for the structural-scaffold reading of \cite{ma2025sketch} that we floated in Section~\ref{sec:gemini}. Re-running \textit{names\_only} at a higher output cap to bound the cap-truncation share of the Gemini accuracy drop is a small but worthwhile follow-up.

\paragraph{What the ablation does not yet cover.} Two cells of the original five-condition matrix in Section~\ref{sec:limitations} remain pending: \textit{regex\_strip} (a non-LLM pleasantry stripper, which would bound what a deterministic non-LLM cleanup achieves) and \textit{llm\_no\_names} (SLM rewrite with the function-name extraction \emph{disabled}, which would test whether the oracle adds anything on top of the rewriter). Given that \textit{names\_only} already falsifies the oracle-only hypothesis, the highest-priority remaining contrast is \textit{llm\_no\_names} vs.\ \textit{full}: a null result there would indicate the oracle can be removed from the production deployment without accuracy cost. Section~\ref{sec:baselines} extends the comparison surface in a different direction: a head-to-head against the LLMLingua-2 \cite{pan2024llmlingua2} prompt compressor at matched rate, which independently reproduces the structural-scaffold destabilisation pattern reported here for \textit{names\_only} on \texttt{gemini-2.5-flash-lite}.

\subsection{Comparison Against Prompt-Compression Baselines}
\label{sec:baselines}

A predictable concern with reporting prompt-side numbers only against the raw cloud call is that any well-tuned post-hoc compressor might in principle recover the same savings without the architectural complexity of an SLM rewrite. We address this directly with a head-to-head comparison against LLMLingua-2 \cite{pan2024llmlingua2}, the strongest publicly available task-agnostic prompt compressor, on the same OMH-Polyglot inputs, the same three backends, and the same three-run protocol at \texttt{temperature}\,$=0$ used in Section~\ref{sec:results}.

\paragraph{Experimental setup.} LLMLingua-2 is a token-classification model (\texttt{xlm-roberta-large}, $\sim\!560$\,M parameters) trained to keep or drop tokens at a target rate; the multilingual XLM-RoBERTa backbone makes it the most natural baseline for a benchmark whose specifications are written in Turkish, Arabic, Chinese, and code-switched mixtures \cite{pan2024llmlingua2}. We run the public checkpoint \texttt{microsoft/llmlingua-2-xlm-roberta-large-meetingbank} through the official \texttt{PromptCompressor} API at \texttt{rate}\,$=0.6$ (a $40\%$ token-reduction target, matched to the prompt-side band achieved by \textit{Ours}), with the recommended \texttt{force\_tokens} list for syntactic punctuation. Because LLMLingua-2 is task-agnostic and is not aware of executable contracts, we detach the \texttt{assert} block from every prompt before compression and concatenate it back verbatim afterwards; without this guard the grader would be silently corrupted by token pruning inside the test lines, which is neither how LLMLingua-2 is meant to be used nor a fair comparison. The Python proxy that hosts LLMLingua-2 is otherwise byte-compatible with the OpenAI-style proxy used for the \textit{Ours} arm: identical system prompt, identical \texttt{<solution>} post-processing, identical $5\%$ token-budget guard. Any per-cell delta is therefore attributable to the compression mechanism alone. The LLMLingua-2 rows of Tables~\ref{tab:main-results} and \ref{tab:dollar-cost} report this matched-rate cell.

\paragraph{Headline contrast.} Across all three backends \textit{Ours} strictly Pareto-dominates LLMLingua-2 on OckScore: $99.08 > 76.91$ on \texttt{gpt-3.5-turbo}, $98.88 > 96.56$ on \texttt{gpt-4o}, and $97.54 > 33.85$ on \texttt{gemini-2.5-flash-lite}. LLMLingua-2 hits the compression rate it was configured for (within-tokenizer prompt reductions of $15$--$21\%$ vs.\ Raw), but the cross-lingual rewrite compresses substantially harder ($34$--$47\%$ vs.\ Raw) and, crucially, preserves accuracy where LLMLingua-2 does not.

\paragraph{Per-backend nuance.} The contrast is most cleanly read by backend strength:
\begin{itemize}
    \item On \texttt{gemini-2.5-flash-lite}, LLMLingua-2 collapses to $35.00\%$ accuracy ($-60.00$\,pp vs.\ Raw, $-63.00$\,pp vs.\ Ours) with a $2.9\times$ answer-token expansion (mean $1{,}026$ tokens/answer vs.\ $359$ in Raw and $327$ in Ours); total tokens \emph{more than double} ($243{,}193$ vs.\ Raw's $116{,}653$) and OckScore drops to $33.85$. The failure mode quantitatively echoes the \textit{names\_only} ablation reported in Section~\ref{sec:ablation} (Gemini at $76.50\%$ accuracy with $1{,}353$ tokens/answer) and is consistent with the structural-scaffold reading we attribute to \cite{ma2025sketch}: the lite model destabilises into runaway generation when the input lacks English semantic scaffolding, regardless of whether the scaffold is removed by bypassing the rewriter (\textit{names\_only}) or by lossy compression of the multilingual body (LLMLingua-2). The two arms therefore independently validate the same mechanism, on the same backend, by two different intervention paths.
    \item On \texttt{gpt-3.5-turbo}, LLMLingua-2 drops accuracy to $77.33\%$ ($-22.17$\,pp vs.\ both Raw and Ours) while saving only $-15.7\%$ on dollar cost vs.\ Ours. The arm sits well below the Pareto frontier (Raw scores OckScore $99.04$ for $\$0.0878$ per $200$-task run; LLMLingua-2 scores $76.91$ for $\$0.0852$): a marginal cost concession with a large accuracy penalty.
    \item On \texttt{gpt-4o}, the strongest backend, LLMLingua-2 is most competitive: $97.17\%$ accuracy ($-2.33$\,pp vs.\ Ours, within $2\times$ the run-level standard deviation reported in Table~\ref{tab:ablation}) at $\$0.9683$ per $200$-task run, an $8.7\%$ dollar saving relative to Ours. The Pareto frontier still favours Ours on OckScore ($98.88$ vs.\ $96.56$), but the contrast is closer here than on the weaker backends---consistent with frontier models recovering from prompt damage that lite models do not.
\end{itemize}

\paragraph{Compression-side latency.} LLMLingua-2 compression on CPU averages $689$\,ms per prompt (median $687$\,ms, $p90$ $703$\,ms) on the same development hardware on which our SLM rewrite averages $176$\,ms per prompt (Section~\ref{sec:results-latency}). The intuition that a $\sim\!560$\,M-parameter token-classification head must be faster than a $3$\,B-parameter SLM is therefore not borne out at this hardware tier: the SLM rewrite is faster per-prompt, in addition to delivering the OckScore and accuracy outcomes above. A CUDA-enabled LLMLingua-2 deployment closes part of this gap but does not change the OckScore ordering.

\paragraph{Offline qualitative evidence of spec damage.} A prerequisite for the live runs above was an offline batch in which LLMLingua-2 compressed every OMH-Polyglot prompt without dispatching to the cloud; this produces a side-by-side raw-vs-compressed corpus that exposes the semantic drift the live runs then quantify in accuracy. Two patterns dominate. First, the \texttt{zh\_pure} style is most damaged: $5/200$ rows lose more than half of their original byte length, including a Roman-numeral specification (\texttt{roman\_to\_int}) whose compressed body is \verb|IV=4,IX=900| where the original gave \verb|IV=4,IX=9,XL=40,XC=90,CD=400,CM=900|---a direct factual corruption, not just terseness. Second, $5/200$ rows lose the target function name from the natural-language body entirely (all in code-switched styles, where the compressor preferentially keeps already-English jargon and drops English identifiers it cannot disambiguate from prose); the asserts still carry the identifier (we preserve them), but the natural-language scaffold the lite model needs is gone. These qualitative observations match the live accuracy ordering: the largest accuracy regression is on \texttt{gemini-2.5-flash-lite}, the backend most sensitive to scaffold loss, exactly as Section~\ref{sec:gemini} predicts.

\paragraph{Scope.} We deliberately limit this comparison to LLMLingua-2 in its as-shipped multilingual configuration. LongCodeZip \cite{shi2025longcodezip}, the alternative code-specific compressor we cite in related work, requires per-token perplexity from a reference LM during compression---the recursive cloud dependency our paper specifically argues against---and is therefore not on the same Pareto axis as our middleware. We also do not extend LLMLingua-2 with per-task function-name protection or other hand-tuned domain hints: doing so would constitute ad-hoc fairness boosts for a baseline we are competing against, and the result reported here is what a practitioner would obtain by using the published model as documented. A protocol that runs LongCodeZip against a local reference LM (e.g.\ the same \texttt{Llama 3.2 3B} we already deploy) would close the remaining published-baseline comparison cleanly; we list it as an immediate follow-up in Section~\ref{sec:limitations}.

\subsection{End-to-End Latency Cost}
\label{sec:results-latency}
Our framework is not free in wall-clock time. Per-prompt local-rewrite latency, computed as $(\bar{t}_{\text{Ours}} - \bar{t}_{\text{Raw}})/200$, is approximately $118$ ms for \texttt{gpt-4o}, $131$ ms for \texttt{gemini-2.5-flash-lite}, and $279$ ms for \texttt{gpt-3.5-turbo}. As a fraction of total benchmark wall-clock, the proxy adds $+10.5\%$ on \texttt{gpt-4o} (where the cloud call dominates), $+40.5\%$ on \texttt{gemini-2.5-flash-lite}, and $+57.5\%$ on \texttt{gpt-3.5-turbo} (where the cheap backend's own latency is small enough that the local rewrite becomes a meaningful fraction of total time). Practitioners should weigh this against the token-cost savings; for interactive IDE use the per-prompt overhead is sub-second on standard development hardware, but it is not zero and we report it explicitly in Section~\ref{sec:limitations}.

\subsection{From Tokens to Dollars: Asymmetric Cost Outcomes}
\label{sec:results-dollars}
A token-only analysis is insufficient for a paper framed around operating cost: commercial APIs price output tokens at $3$--$4\times$ the input rate, so the same $-8\%$ change in total tokens can translate into very different dollar outcomes depending on which side of the prompt the savings come from. Table~\ref{tab:dollar-cost} reports the full cost-per-$200$-task amount for each backend using publicly listed pricing as of May~$2026$\footnote{$\$0.50/\$1.50$ per million input/output tokens for \texttt{gpt-3.5-turbo}; $\$2.50/\$10.00$ for \texttt{gpt-4o}; $\$0.10/\$0.40$ for \texttt{gemini-2.5-flash-lite}.}.

\begin{table}[h]
\centering
\small
\setlength{\tabcolsep}{4pt}
\begin{tabular}{lcccrr}
\toprule
\textbf{Cloud Agent} & \textbf{Raw} & \textbf{LLMLingua-2} & \textbf{Ours} & \textbf{$\Delta\$$ LL2} & \textbf{$\Delta\$$ Ours} \\
\midrule
\texttt{gpt-3.5-turbo}        & \$0.0878 & \$0.0852 & \$0.1011 & $-2.9\%$  & $+15.1\%$ \\
\texttt{gpt-4o}               & \$1.0641 & \$0.9683 & \$1.0601 & $-9.0\%$  & $-0.4\%$ \\
\texttt{gemini-2.5-flash-lite}& \$0.0332 & \$0.0859 & \$0.0291 & $+158.7\%$ & $-12.4\%$ \\
\bottomrule
\end{tabular}
\caption{Per-200-task cloud-API cost in USD, using listed input/output prices as of May 2026. Both $\Delta\$$ columns are relative to that backend's Raw cost. Local-compute cost (Llama 3.2 3B at $\sim$$100$\,W for $\sim$$176$\,ms per query) is approximately \$0.07 per $100{,}000$ queries and is omitted from the table as it is dominated by every cell shown; LLMLingua-2's compression cost on CPU ($\sim$$689$\,ms per query, Section~\ref{sec:baselines}) is similarly negligible in absolute dollar terms and omitted on the same basis. Note that on \texttt{gpt-4o} the LLMLingua-2 baseline saves money relative to \textit{Ours} ($\$0.9683$ vs.\ $\$1.0601$) but at a $-2.33$\,pp accuracy cost; the OckScore ordering in Table~\ref{tab:main-results} still favours \textit{Ours} ($98.88$ vs.\ $96.56$).}
\label{tab:dollar-cost}
\end{table}

The dollar column tells a more nuanced story than the token column. On \texttt{gemini-2.5-flash-lite}, where our middleware reduces both prompt and answer tokens, savings are real and meaningful ($-12.4\%$). On \texttt{gpt-4o}, the $-34.0\%$ prompt-token reduction is almost exactly cancelled by the $+3.5\%$ answer-token increase under that model's $4\times$ output-price multiplier, leaving the bill essentially flat. On \texttt{gpt-3.5-turbo}, the proxy actually \textit{increases} the bill by $15.1\%$: the model's $42.3\%$ longer completions are charged at the $3\times$ output rate, more than erasing the input-side savings.

The honest synthesis is that our token-economic mechanism is most valuable when (i) the backend has a small output/input price ratio or (ii) the answer-side does not regress under the cleaner prompt. Both conditions hold for \texttt{gemini-2.5-flash-lite} and the savings compound; only one holds for \texttt{gpt-4o} (flat); neither holds for \texttt{gpt-3.5-turbo} (regresses). We treat this as an additional reason to prefer the middleware for cost-sensitive deployments on lite/efficient backends, which is precisely the segment where token-economic optimisation matters most. We also acknowledge it as a real limitation on the headline claim that the system universally reduces operating cost.

The LLMLingua-2 baseline column in Table~\ref{tab:dollar-cost} adds one further nuance worth surfacing here. On \texttt{gpt-4o} a token-classification compressor at matched rate \emph{does} beat \textit{Ours} on dollar cost ($-9.0\%$ vs.\ Raw, vs.\ $-0.4\%$ for \textit{Ours}), because the strongest backend recovers gracefully from prompt damage and pays the smaller answer bill. That dollar saving comes with a $-2.33$\,pp accuracy concession, however, so the OckScore ordering in Table~\ref{tab:main-results} still places \textit{Ours} ahead ($98.88$ vs.\ $96.56$). On the other two backends LLMLingua-2 is strictly worse: marginal cost saving with catastrophic accuracy loss on \texttt{gpt-3.5-turbo} ($-22.17$\,pp for $-2.9\%$ dollars), and outright cost regression with severe accuracy loss on \texttt{gemini-2.5-flash-lite} ($-63.00$\,pp accuracy and $+158.7\%$ dollar cost vs.\ Raw, attributable to the $2.9\times$ answer-token explosion documented in Section~\ref{sec:baselines}). The Pareto-aware reading is that token-classification compression buys some dollar headroom on output-expensive frontier backends, but only at the cost of accuracy our cross-lingual rewriter does not pay; on output-cheap lite backends it is worse on every axis.

\section{Conclusion}
\label{sec:conclusion}
We have presented a localized, cost-arbitrage middleware that intervenes between the developer and the cloud agent, performs cross-lingual translation and structural rewriting on edge hardware, and forwards a hyper-dense English specification to the cloud backend. On \textit{OMH-Polyglot}---the polyglot-specification member of our \textit{OMH} (\textbf{O}ck-\textbf{M}ultilingual-\textbf{H}ard) benchmark family (Section~\ref{sec:omh-family}), released to surface tokenization overhead that the paired wrapper-only control \textit{OMH-Wrapped} does not induce---the middleware delivers a uniform $34$--$47\%$ reduction in prompt tokens across \texttt{gpt-4o}, \texttt{gpt-3.5-turbo}, and \texttt{gemini-2.5-flash-lite}, and total-token reductions of $8.3$--$18.8\%$ over three independent runs per cell. Accuracy is held at $99.5\%$ on \texttt{gpt-3.5-turbo}, increases nominally by $+1.17$\,pp on \texttt{gpt-4o} (within run-level noise), and increases by $+3.00$\,pp on \texttt{gemini-2.5-flash-lite}.

To attribute the accuracy column to a specific mechanism, we run a \textit{names\_only} ablation that retains the deterministic function-name oracle but drops the SLM rewrite. The oracle in isolation \emph{decreases} accuracy by $20.50$\,pp on \texttt{gpt-3.5-turbo} and $18.50$\,pp on \texttt{gemini-2.5-flash-lite}, and is neutral on \texttt{gpt-4o}, falsifying the hypothesis that identifier salience explains the \textit{Ours} arm's accuracy gains. The SLM rewrite, not the oracle, is the active mechanism; on \texttt{gemini-2.5-flash-lite} it additionally prevents a $3.8\times$ answer-token explosion that the oracle-only arm exhibits. We retain the cluster-level caveat from Section~\ref{sec:metrics} for the precise magnitude of the $+3.00$\,pp Gemini gain (six row-flips concentrated within $\le 6$ algorithm clusters), but the qualitative attribution is no longer in question.

We are deliberate about what the data does \emph{not} show. Answer-side compression is not universal: only \texttt{gemini-2.5-flash-lite} produces shorter completions under our full rewrite, while OpenAI models produce slightly to substantially longer ones, more than offset by prompt-side savings. The previously reported ``Optimization Floor'' on \texttt{gpt-3.5-turbo} disappears under the cleaner experimental protocol and we now treat that earlier finding as an artifact---though the \textit{names\_only} ablation shows the model can still be cratered to $79\%$ accuracy if the multilingual body is forwarded without the SLM rewrite, which reframes the ``floor'' as a symptom of un-translated input rather than a model-capability ceiling.

Translating these token deltas into dollar deltas at listed cloud prices reveals a finding the token-only view conceals: cost outcomes are asymmetric across backends. The middleware saves $12.4\%$ on \texttt{gemini-2.5-flash-lite} but is flat on \texttt{gpt-4o} and \textit{regresses} by $15.1\%$ on \texttt{gpt-3.5-turbo}, where the $3\times$ output/input price ratio combined with longer post-rewrite completions inverts the token-side win. We therefore recommend our middleware for output-cheap or output-stable backends and present this asymmetry honestly rather than averaging it away.

A head-to-head comparison against LLMLingua-2 \cite{pan2024llmlingua2} at matched compression rate (Section~\ref{sec:baselines}) shows our middleware strictly Pareto-dominates token-classification compression on OckScore across all three backends ($99.08 > 76.91$ on \texttt{gpt-3.5-turbo}, $98.88 > 96.56$ on \texttt{gpt-4o}, $97.54 > 33.85$ on \texttt{gemini-2.5-flash-lite}). LLMLingua-2's accuracy regression on the two weaker backends ($-22.17$ and $-63.00$\,pp respectively), accompanied by a $2.9\times$ answer-token expansion on \texttt{gemini-2.5-flash-lite}, mirrors the destabilisation pattern observed in our \textit{names\_only} ablation; the two arms therefore independently validate the structural-scaffold attribution we make for the Gemini accuracy gain in Section~\ref{sec:gemini}. We treat this Pareto-dominance result as the empirical answer to the question ``does the cross-lingual rewrite mechanism actually outperform task-agnostic prompt compression on a multilingual coding benchmark?'': it does, by every margin the OckScore framework can measure.

By programmatically using English as a token-efficient pivot language, our middleware also reduces the linguistic barrier for non-English-speaking developers without forcing them to write English at the IDE surface \cite{ahuja2025efficientxlang}. Future work will integrate (i) the two remaining cells of the ablation matrix (\textit{regex\_strip}, \textit{llm\_no\_names}), (ii) an algorithm-cluster bootstrap with Holm-corrected significance over the headline contrasts, (iii) a comparison against LongCodeZip \cite{shi2025longcodezip} running on a local reference LM (the only remaining published baseline our paper's no-recursive-cloud-cost stance permits), and (iv) repository-aware grounding of the \texttt{[CONTEXT]} block from static analysis output, so that the Tri-Block can replace the Bi-Block deployed in the present benchmark harness.

\section{Limitations}
\label{sec:limitations}
The limitations below are ordered from those that most strongly threaten the validity of our accuracy claims to those that affect deployment scope.

\paragraph{Function-name oracle: confound now resolved against the oracle hypothesis.} The benchmark proxy extracts target function identifiers from the executable \texttt{assert} lines via a deterministic regex (\texttt{functionNamesFromAsserts}) and either instructs the SLM to reproduce them verbatim or, if the SLM omits them, prepends them to the rewritten prompt. The cloud model on the \textit{Ours} arm therefore always receives the target identifier at the top of the \texttt{[TASK]} block, which raised a legitimate concern that the accuracy deltas in Table~\ref{tab:main-results} could be driven by identifier salience rather than by the rewriting mechanism. We address this directly with the \textit{names\_only} ablation reported in Section~\ref{sec:ablation}: the oracle in isolation \emph{decreases} accuracy by $20.50$\,pp on \texttt{gpt-3.5-turbo} and $18.50$\,pp on \texttt{gemini-2.5-flash-lite}, and is neutral on \texttt{gpt-4o} ($+0.17$\,pp, within run-level noise). The oracle therefore cannot explain the \textit{Ours} arm's accuracy results---if it could, \textit{names\_only} would approximate \textit{full}, but it falls $20.50$, $1.00$, and $21.50$\,pp short on the three backends. The remaining limitation in this category is that we have not yet run \textit{llm\_no\_names} (the rewriter without the oracle) to test whether the oracle adds anything on top of the rewriter; we expect a null effect there, but a non-null result would change the recommended deployment configuration.

\paragraph{Pending ablation cells.} The five-condition matrix originally proposed for full attribution is:
\begin{enumerate}
    \item \textit{raw} --- direct cloud call, no proxy. \textbf{Done} (Table~\ref{tab:main-results}).
    \item \textit{names\_only} --- raw prompt with the regex-extracted function name prepended; no LLM rewrite. \textbf{Done} (Table~\ref{tab:ablation}).
    \item \textit{regex\_strip} --- regex-based pleasantry stripper without an LLM. \textit{Pending.}
    \item \textit{llm\_no\_names} --- SLM rewrite with function-name extraction disabled. \textit{Pending.}
    \item \textit{full} --- SLM rewrite with function-name extraction (the \textit{Ours} arm). \textbf{Done} (Table~\ref{tab:main-results}).
\end{enumerate}
The two pending cells are not required to defend the headline accuracy attribution, which is already established by the \textit{raw}/\textit{names\_only}/\textit{full} contrast in Section~\ref{sec:ablation}. They would, however, strengthen the deployment story (\textit{regex\_strip} bounds the non-LLM cleanup floor; \textit{llm\_no\_names} tests whether the oracle is necessary alongside the rewriter).

\paragraph{Descriptive rather than inferential statistics.} OMH-Polyglot is structured as $20$ algorithmic clusters of $10$ multilingual style transforms with shared reference solutions and shared \texttt{assert} suites per cluster, so the effective independent-trial count for accuracy is approximately $20$ rather than $200$; a defensible standard error must therefore be computed at the cluster level (e.g.\ a per-algorithm bootstrap that resamples clusters with replacement and the styles within them as a unit). The three repeats per cell at $\texttt{temperature}{=}0$ characterise residual cross-region API non-determinism rather than population variability. The small accuracy deltas (the $+1.17$\,pp \texttt{gpt-4o} contrast in Table~\ref{tab:main-results}, and the $+0.17$\,pp \textit{names\_only}-vs-\textit{raw} contrast on the same backend in Table~\ref{tab:ablation}) fall within the run-level standard deviation and are reported as descriptive only. The large ablation deltas in Table~\ref{tab:ablation} ($-18.50$ to $-20.50$\,pp on the two weaker backends) are well outside any plausible cluster-level noise band and should survive a clustered analysis; we still mark them as pending such an analysis for completeness, and a Holm-corrected significance pass over the $\{$3 backends $\times$ $\{$\textit{full}--\textit{raw}, \textit{names\_only}--\textit{raw}$\}\}$ contrast set is a required follow-up.

\paragraph{End-to-end latency overhead.} The proxy adds $118$--$279$\,ms per prompt depending on the backend (Section~\ref{sec:results-latency}), raising total benchmark wall-clock by $+10.5\%$ on \texttt{gpt-4o}, $+40.5\%$ on \texttt{gemini-2.5-flash-lite}, and $+57.5\%$ on \texttt{gpt-3.5-turbo}. Token cost is reduced at the expense of wall-clock cost; deployment in latency-bound interactive settings must weigh this trade-off explicitly.

\paragraph{Backend-specific dollar regression.} The dollar-cost outcome is asymmetric (Section~\ref{sec:results-dollars}): the middleware saves $-12.4\%$ on \texttt{gemini-2.5-flash-lite}, is flat ($-0.4\%$) on \texttt{gpt-4o}, and increases cost by $+15.1\%$ on \texttt{gpt-3.5-turbo} because that backend prices output tokens at $3\times$ input and the cleaner prompt elicits longer completions. Deployment recommendations are therefore conditional on the backend's output/input price ratio and on prior knowledge of its answer-token behaviour under structured prompts.

\paragraph{Translation quality.} OMH-Polyglot's Turkish segments are native-quality. The Modern Standard Arabic and Simplified Chinese segments use conservative technical register and have not yet been audited by native speakers; an audit pass is scheduled before any external release of the dataset.

\paragraph{Schema mismatch between benchmark and IDE.} The benchmark proxy emits a Bi-Block (\texttt{[CONTEXT]}, \texttt{[TASK]}); the IDE extension extends this with a third \texttt{[CONSTRAINTS]} block for workspace-level constraints not encoded as test code. The results in Section~\ref{sec:results} reflect the Bi-Block configuration only.

\paragraph{Outstanding code-specific baseline.} Section~\ref{sec:baselines} reports a head-to-head comparison against LLMLingua-2 \cite{pan2024llmlingua2} at matched compression rate; \textit{Ours} strictly Pareto-dominates on OckScore across all three backends, with LLMLingua-2's accuracy regression most severe on the lite backend ($-63.00$\,pp) and mildest on the frontier backend ($-2.33$\,pp). The remaining published code-specific compressor we cite, LongCodeZip \cite{shi2025longcodezip}, requires per-token perplexity from a reference LM during compression---the recursive cloud dependency our paper specifically argues against---and is therefore not directly comparable on our Pareto axis. A protocol that runs LongCodeZip against a local reference LM (e.g.\ the same \texttt{Llama 3.2 3B} we already deploy) would close this comparison cleanly and is listed as an immediate follow-up in Section~\ref{sec:conclusion}. We also acknowledge that the LLMLingua-2 cell is reported at a single matched rate ($r{=}0.6$); a small Pareto sweep at $r \in \{0.5, 0.7\}$ would tighten the contrast against a hypothetically harder-tuned baseline, but does not affect the OckScore ordering given the magnitude of the lite-backend accuracy gap.

\paragraph{Hardware footprint.} Local execution of \texttt{Llama 3.2} (3B) requires approximately $3.4$\,GB of dedicated memory to sustain sub-second inference latency, which may preclude deployment on memory-constrained or legacy hardware \cite{meta2024llama32}.

\paragraph{Syntactic validation.} The post-rewrite validator (Section~\ref{sec:guard}) performs regex-level legality checks (block markers present, no leaked code block, no filled \texttt{<solution>} envelope, minimum length) and a $5\%$ character-based token-budget guard, with up to two repair attempts and a raw-prompt fallback. It does not extract the requirement set from the raw query, does not compute requirement coverage, and does not detect identifier drift beyond the function-name hint. The guard ensures the cloud-billed payload is never inflated relative to the raw baseline, but does not guarantee semantic completeness; replacing this with a requirement-coverage pass is on the follow-up roadmap.

To facilitate transparency and reproducibility, the IDE extension, the benchmark proxy, the OMH-Polyglot dataset, and the analysis scripts are released in the supplementary materials.\footnote{Anonymized repository link: \url{https://anonymous.4open.science/r/example-id}}

\bibliographystyle{plain}   % or ieeetr, unsrt, etc.
\bibliography{citations}

\end{document}